\section{Fazit und Ausblick} \label{sec:fazit}

\subsection{Zusammenfassung der Ergebnisse} \label{subsec:zusammenfassung}
Die vorliegende Arbeit untersuchte die Realisierbarkeit eines hybriden Systems zur Kontrolle von Zugangs- und Auslassberechtigungen für Schulen unter Verwendung von RFID-Technologie und lokaler KI-Inferenz. Im Fokus stand die Beantwortung der Leitfrage: \textit{„Wie lässt sich die Kontrolle von Zugangs- und Auslassberechtigungen an Schulen durch die Kombination von passiver RFID-Technologie und lokaler KI-Gesichtserkennung effizient, fälschungssicher sowie datenschutzkonform digitalisieren, um papierbasierte Prozesse abzulösen?“}

Die Evaluation der Benchmarks auf einem Raspberry Pi 4 zeigte, dass die technische Umsetzung die gesteckten Ziele erreicht. Die Hypothese \textbf{H1} wurde verifiziert: Die RFID-Validierung erfolgt sofort (ca. 8\,ms), und auch der optionale biometrische Verifikationsprozess bleibt mit ca. 3,5 Sekunden in einem praxistauglichen Rahmen. Hypothese \textbf{H2} zur Zuverlässigkeit wurde bestätigt (100\,\% Erkennungsgenauigkeit im Test, Zielwert: 90\,\%). Schließlich wurde die Hypothese \textbf{H3} plausibilisiert: Durch den Einsatz von Edge-Computing auf einem lokalen Raspberry Pi verlassen biometrische Daten das Schulnetzwerk nie, was „Privacy by Design“ gewährleistet.

\subsection{Beantwortung der Leitfrage} \label{subsec:beantwortung_leitfrage}
Die Leitfrage lässt sich wie folgt beantworten: Eine effiziente, fälschungssichere und datenschutzkonforme Zugangs- und Auslasskontrolle lässt sich durch eine dezentrale Full-Stack-Architektur realisieren, welche biometrische Inferenz am „Edge“ als zweiten Faktor ausführt.

\textbf{Sicherheit} wird durch die Kombination zweier Faktoren erreicht: Der RFID-Kartenbesitz wird durch biometrische Identitätsprüfung ergänzt, was Kartenmissbrauch verhindert. Die Bildnormalisierung (HEIC-zu-PNG) stellt Kompatibilität sicher.

Die \textbf{Datenschutzkonformität} wird durch lokale Verarbeitung sichergestellt. Im Gegensatz zu Cloud-Lösungen entfallen riskante Datenübertragungen. Die Nutzung von SQLite minimiert die Angriffsfläche für externe Zugriffe.

Die \textbf{Vollautomatisierung} gelingt durch das Backend, das Berechtigungen rollenbasiert (RBAC) verwaltet und wiederkehrende Auslassmuster erkennt. Das System entlastet das Personal und ermöglicht Lehrkräften die granulare Verwaltung von Ausnahmegenehmigungen per App.

\subsection{Wissenschaftlicher und praktischer Mehrwert} \label{subsec:mehrwert}
Der wissenschaftliche Mehrwert liegt in der Demonstration, dass leistungsfähige KI-Modelle zur Gesichtserkennung auch auf Consumer-Hardware (Raspberry Pi) performant und datenschutzkonform eingesetzt werden können. Praktisch bietet das System eine skalierbare und kosteneffiziente Lösung für Bildungseinrichtungen.

\subsection{Ausblick und zukünftige Arbeiten} \label{subsec:ausblick}
\begin{enumerate}
    \item \textbf{Echtzeit-Videoerkennung:} Die Migration von statischen Schnappschüssen zu einer kontinuierlichen Videoanalyse (Live-Stream) könnte den Prozess weiter beschleunigen.
    \item \textbf{Anti-Spoofing-Maßnahmen:} Um die Sicherheit zu erhöhen, sollten Mechanismen zur Lebenderkennung (Liveness Detection) integriert werden. Durch die Analyse feiner Veränderungen im Gesicht (Mikrobewegungen, Blinzeln) im Videostream lässt sich das Überlisten durch vorgehaltene Fotos verhindern.
    \item \textbf{PostgreSQL-Migration:} Für sehr große Schulen ist eine Migration von SQLite zu einer PostgreSQL-Datenbank sinnvoll, um die Schreibperformanz zu optimieren.
    \item \textbf{Integration in Schulverwaltungssysteme:} Eine Anbindung an Stundenplansoftware (z.\,B. Untis, EduPage) sowie Entschuldigungs- oder Krankmeldungssysteme würde die automatisierte Vergabe von Berechtigungen in Freistunden, nach Schulschluss oder bei Krankheit ermöglichen.
\end{enumerate}

Abschließend lässt sich festhalten, dass das Projekt „Digital School Access Control“ ein solides Fundament für eine moderne Schul-Sicherheitsarchitektur bildet.


